\chapter{Cut Elimination and Analytic Derivations}
\label{ch:cut-elimination}

\begin{hbdefinition}{def:p07-analytic-derivation}{Analytic derivations in the displayed calculi}
In this chapter, a propositional \(\LK\)- or \(\LJ\)-derivation is
\emph{analytic} when it is cut-free and every formula occurrence in it is
a subformula of a formula in its endsequent.  Subformulas are computed in
the selected calculus's primitive syntax, after expanding its declared
abbreviations.

For the first-order versions, an analytic derivation \emph{modulo terms}
is cut-free and has the following property: every formula occurrence is,
up to bound-variable renaming, a capture-avoiding term-substitution instance
of a subformula of an endsequent formula.  A subformula is considered on
its own when applying the substitution, so a variable bound by a removed
outer quantifier is free in that subformula.  Terms and eigenparameters
must belong to the language permitted by the selected quantifier rules,
and each inference still satisfies its own freshness conditions.

For example, the premise formula \(P(t)\) of an instantiation of
\(\forall x\,P(x)\) is an instance of its subformula \(P(x)\),
although \(P(t)\) need not itself occur literally in the quantified
formula.  This is the term qualification used below; it can admit
infinitely many instances and therefore does not make first-order proof
search finite.  The subsequent subformula lemma and cut-elimination
argument establish the corresponding existence claims at their stated
proof status.  Extensions of these calculi require their own such results.
\index{analytic derivation}
\index{subformula property!modulo terms}
\end{hbdefinition}

\section{Measures on cuts}
\label{sec:cut-elimination-01}

\begin{definition}[Formula rank and derivation height]
The rank \(r(A)\) of a formula is the number of logical symbols in it.  The
height \(h(D)\) of a derivation is the greatest number of inference steps on a
branch from a leaf to its root.  These measures are different: formula rank
measures syntax, whereas derivation height measures a proof tree.
\end{definition}

The reduction below is run for the following generalization of cut, the
\emph{mix} (multicut) rule.  Here \(A^{k}\) abbreviates \(k\)
comma-separated occurrences of \(A\).  For \(\LK\), with \(m,n\geq1\),
and with \(\Delta\) and \(\Sigma\) permitted to contain further
occurrences of \(A\):
\[
\frac{\Gamma\Rightarrow\Delta,A^{m}
      \qquad
      A^{n},\Pi\Rightarrow\Sigma}
     {\Gamma,\Pi\Rightarrow\Delta,\Sigma}\;\textsc{Mix}.
\]
For \(\LJ\) the succedent is single, so the left premise's succedent
multiplicity \(m\) is fixed at one: from \(\Gamma\Rightarrow A\) and
\(A^{n},\Pi\Rightarrow C\) with \(n\geq1\), \textsc{Mix} concludes
\(\Gamma,\Pi\Rightarrow C\).  The displayed occurrences of \(A\) are the
\emph{mix formula} of the instance.
